% =====================================================================
\section{Application Integration and Live Captioning}
\label{sec:application}
% =====================================================================

\subsection{Batch prototype versus streaming application}
Two kinds of application are possible, and the distinction determines how much
engineering the demonstration needs.
\begin{itemize}
  \item A \textbf{batch prototype} processes a complete recorded clip---uploaded
        or recorded in the browser---and returns its caption. It reuses the
        evaluation code almost unchanged and its correctness is easy to verify.
  \item A \textbf{streaming application} captions speech while it is being spoken.
        It needs buffering, synchronisation of two live devices, and repeated
        inference on overlapping windows. Because the encoders are bidirectional,
        the model needs future context and cannot emit a character the moment a
        frame arrives; streaming is therefore approximated by re-decoding recent
        windows, which introduces latency of roughly one window.
\end{itemize}

The recommended plan \rec\ is to build the batch prototype first (Mode~A) and
then a windowed pseudo-streaming mode (Mode~B) that reuses its preprocessing and
inference functions. Mode~A is the guaranteed deliverable; Mode~B is the live
captioning demonstration. A fully streaming model with unidirectional encoders
is outside the six-week scope.

\subsection{Mode A: batch captioning in Gradio}
A Gradio page offers a webcam-and-microphone recorder and a file-upload box. On
submission, the server extracts frames with OpenCV and audio with
\texttt{ffmpeg}, resamples video to 25\,fps and audio to 16\,kHz, runs the visual
and audio pipelines, the model and beam search, and displays the caption beside
the video. The page also shows the mouth-crop sequence and a plot of the
per-frame fusion gate, which makes the system's behaviour visible during the
presentation. An optional control adds MUSAN noise at a chosen SNR to the clip,
so that the audience can see captions with and without noise on the same
recording.

\subsection{Mode B: windowed live captioning}

\begin{figure}[H]
\centering
\begin{tikzpicture}[
  node distance=6mm and 8mm,
  th/.style={draw=lsblue, rounded corners=2pt, fill=lsblue!6, align=center,
             font=\scriptsize\sffamily, minimum height=11mm, text width=27mm},
  buf/.style={draw=black!50, fill=black!4, align=center, font=\scriptsize\sffamily,
              minimum height=8mm, text width=27mm},
  ar/.style={-{Stealth[length=2mm]}, thick}]
\node[th] (cam) {Camera thread\\OpenCV, timestamped frames};
\node[th, below=of cam] (mic) {Audio callback\\sounddevice, 16\,kHz blocks};
\node[buf, right=of cam] (vb) {Frame ring buffer\\(last 4\,s)};
\node[buf, right=of mic] (ab) {Audio ring buffer\\(last 4\,s)};
\node[th, right=of $(vb)!0.5!(ab)$, text width=30mm] (inf) {Inference worker\\every 1\,s: extract 3\,s window,\\crop, features, model, beam search};
\node[th, right=of inf, text width=24mm] (ui) {Display\\partial / final caption overlay};
\draw[ar] (cam) -- (vb); \draw[ar] (mic) -- (ab);
\draw[ar] (vb) -- (inf); \draw[ar] (ab) -- (inf); \draw[ar] (inf) -- (ui);
\end{tikzpicture}
\caption{Threads and buffers of the windowed live-captioning mode.}
\label{fig:live}
\end{figure}

\paragraph{Input and buffering.} A camera thread reads frames with OpenCV and
stores each with a timestamp from a single monotonic clock. A
\texttt{sounddevice} callback appends 16\,kHz audio blocks, also timestamped, to a
ring buffer. Both buffers hold the most recent four seconds, which bounds memory
use (four seconds of $640\times480$ colour frames at 30\,fps is about 110\,MB; the
buffer can store downscaled frames or pre-computed crops to reduce this).

\paragraph{Temporal synchronisation.} Webcams rarely deliver exactly 25\,fps and
may drop frames. For a window covering times $[t_0, t_0+3\,\text{s})$ the worker
selects, for each of the 75 target instants $t_0 + 40k$\,ms, the buffered frame
with the nearest timestamp, and cuts the matching 48{,}000 audio samples. Any
constant offset between the camera and microphone pipelines is measured once
with a hand-clap test (comparing the frame where the hands meet with the audio
onset) and subtracted \rec.

\paragraph{Windowing and preprocessing.} Every second the worker processes the
latest 3\,s window, matching the duration of GRID training utterances and giving
$T=75$, which satisfies the CTC length condition (\cref{eq:ctc-cond}). Face Mesh
runs on each new frame as it arrives rather than once per window, so each frame is
processed only once even though windows overlap; the crop smoothing uses an
exponential moving average over past frames. Audio features are computed with
the stored training statistics.

\paragraph{Inference and decoding.} The model runs in evaluation mode, with
gradient tracking disabled by \texttt{torch.inference\_mode()}, and its output
is decoded by beam search. If more than half of
the window's frames lack a face, the visual stream is replaced by
$\mathbf{e}_{\text{miss}}$ and the caption is marked ``audio only''.

\paragraph{Partial and final captions.} A simple energy-based voice-activity
detector tracks speech on the audio buffer. While speech continues, the caption
for the latest window is shown as a grey \emph{partial} caption and replaced each
second. When the detector reports at least 500\,ms of silence, the worker decodes
one window aligned to the detected utterance and displays it as a \emph{final}
caption in black, appended to a short caption history. Restricting final
captions to detected utterances avoids the difficult problem of merging text
from overlapping windows.

\paragraph{Errors and missing input.} If the camera or microphone cannot be
opened, the application reports it and continues in the remaining modality
where possible. Exceptions in the inference worker are logged and the window is
skipped, so a single failure does not stop the display loop.

\subsection{Latency, cost and memory}
Caption latency in Mode~B is the sum of waiting for the window to fill (up to
1\,s for partial captions, the end-of-speech delay for final ones), Face Mesh on
the latest frames, and model and beam-search time. The 3D convolutions over 75
frames dominate model cost; the per-window cost is the same as one test
utterance and will be measured as described in \cref{sec:evaluation}. Real-time
audio-visual recognition on modest hardware is still an active research
problem~\cite{ryumin2023,aripin2024}, so these measurements, rather than
assumptions, will decide the final window settings. If latency
is too high, the team will, in order: run inference every 1.5\,s instead of 1\,s,
reduce beam width, run the model on a GPU, or quantise the model weights.

\subsection{Testing the application}
\begin{itemize}
  \item \textbf{Pre-recorded clips.} Held-out GRID test clips are played through
        Mode~A and through Mode~B (by feeding the file to the buffers in real time).
        Their captions must match those produced by the evaluation script for the
        same clips, which confirms that the application path is equivalent to the
        evaluated path.
  \item \textbf{Live input.} Team members read GRID-grammar sentences to the
        webcam in a quiet room and with played-back background noise. Because the
        model is trained on GRID, it is expected to work only on GRID-like sentences
        spoken in similar framing; free conversation is outside its training
        distribution, and the demonstration will say so.
  \item \textbf{Failure cases.} Covering the mouth, turning the head and leaving the
        frame must produce ``audio only'' captions rather than crashes.
\end{itemize}
